\subsection{Chips for the base instruction set}
\label{sec:base}

For each instruction family we give the chip, its frame, how the row fixes the result, and the lookups the chip adds to its frame's. Operands follow Definition~\ref{def:format}: $a$ is the word written to (or, for stores and branches, read from) register $\mathsf{op\_a}$, and $b$ and $c$ are the words of $\mathsf{op\_b}$ and $\mathsf{op\_c}$, from registers or decoded immediates, in the notation of \S\ref{sec:air-chips}. Every result is gated by $(1-\mathsf{op\_a\_0})$, which the tables omit.

\paragraph{Frames.} The frame's interactions are paid on every row (\S\ref{sec:air-frame}): 20 for the R-type frame and 16 for the I-type and shift-amount frames, itemised in Appendix~\ref{sec:app-frame}. The \emph{universal} frame is the R-type frame with its reads gated by the immediate flags. The four frames play the part of OpenVM's six RV32IM adapters~\citep{openvm}; loads, stores and branches share the I-type frame. Widths are from Table~\ref{tab:chips}, and lookups are counted beyond the frame's.

\subsubsection{Logical instructions}
\label{sec:base-logical}

\begin{center}
\small
\begin{tabular}{@{}llll@{}}
\toprule
Instructions & Chip (width), frame & Result & Added lookups \\
\midrule
\texttt{and}, \texttt{or}, \texttt{xor}, \texttt{nor} & \texttt{Bitwise} (37), R & $a_i = \mathrm{OP}(b_i, c_i)$, $i < 4$ & 4 byte-table \\
\texttt{andi}, \texttt{ori}, \texttt{xori} & \texttt{BitwiseImm} (33), I & $a_i = \mathrm{OP}(b_i, c_i)$, $i < 4$ & 4 byte-table \\
\bottomrule
\end{tabular}
\end{center}

A bitwise operation has no constraint at all: each byte pair is one lookup into the byte table, whose operation index is a selector-weighted sum of the opcode flags (\S\ref{sec:air-bytes}). The immediates of \texttt{andi}, \texttt{ori} and \texttt{xori} are zero-extended at decode, and \texttt{nor} has no immediate form.

\subsubsection{Arithmetic instructions}
\label{sec:base-arith}

\begin{center}
\small
\begin{tabular}{@{}llll@{}}
\toprule
Instructions & Chip (width), frame & Result & Added lookups \\
\midrule
\texttt{add}, \texttt{addu} & \texttt{AddSub} (36), R & $a = b + c \bmod 2^{32}$ & none \\
\texttt{sub}, \texttt{subu} & \texttt{AddSub} (36), R & $a + c = b \bmod 2^{32}$ & none \\
\texttt{addi}, \texttt{addiu} & \texttt{AddSubImm} (33), I & $a = b + \mathrm{sext}(\mathit{imm}) \bmod 2^{32}$ & none \\
\texttt{lui} & \texttt{AddSubImm} (33), I & $a = 0 + (\mathit{imm} \ll 16)$ & none \\
\texttt{mfhi}, \texttt{mflo}, \texttt{mthi}, \texttt{mtlo} & \texttt{AddSubImm} (33), I & $a = b + 0$ & none \\
\bottomrule
\end{tabular}
\end{center}

Addition is checked byte by byte with carries that are not columns: the carry out of byte $i$ is the linear expression $\theta_i = (b_i + c_i - a_i + \theta_{i-1})/256$, $\theta_{-1} = 0$, and the chip asserts $\theta_i(\theta_i - 1) = 0$; subtraction asserts the same of $a + c = b$. The frame byte-checks the result, so the family adds four columns (two selectors, two gates) and no lookup. Jolt instead forms $b + c$ as one element of a 256-bit field and looks up its low bits~\citep{jolt}, which $2^{32} > p$ rules out here. \texttt{lui} adds an immediate shifted at decode to register 0, and the moves to and from \texttt{HI} and \texttt{LO} add zero (\S\ref{sec:overview-format}). \texttt{add} and \texttt{addi} decode as \texttt{addu} and \texttt{addiu} (\S\ref{sec:isa-emulator}).

\subsubsection{Comparisons}
\label{sec:base-slt}

\begin{center}
\small
\begin{tabular}{@{}llll@{}}
\toprule
Instructions & Chip (width), frame & Result & Added lookups \\
\midrule
\texttt{sltu}, \texttt{slt} & \texttt{Lt} (50), R & $a = x_s(1-y_s) + \mathsf{eq}(x_s, y_s)\,\mathsf{LTU}(x_{<s}, y_{<s})$ & 2 \texttt{AND}, 1 \texttt{LTU} \\
\texttt{sltiu}, \texttt{slti} & \texttt{LtImm} (47), I & as above, $y = \mathrm{sext}(\mathit{imm})$ & 2 \texttt{AND}, 1 \texttt{LTU} \\
\bottomrule
\end{tabular}
\end{center}

With $x = b$ and $y = c$ the result is Jolt's signed less-than, while for the unsigned forms the sign bits are not split off and \texttt{LTU} compares the full words~\citep{jolt}. Two \texttt{AND} lookups clear the top bit of each high byte, which gives $x_{<s}$, $y_{<s}$ and, by difference, the sign bits. The unsigned part is decomposed differently: Jolt sums over chunks, weighting each chunk's less-than by the equality of every chunk above it, a less-than and an equality lookup per chunk; \texttt{Lt} witnesses a one-hot flag for the most significant differing byte, asserts by constraints that the bytes above it agree and, through a witnessed inverse, that the flagged bytes differ, and spends one \texttt{LTU} lookup on the flagged pair. The upper three result bytes are zero.

\subsubsection{Shifts and rotates}
\label{sec:base-shift}

\begin{center}
\small
\begin{tabular}{@{}llll@{}}
\toprule
Instructions & Chip (width), frame & Result & Added lookups \\
\midrule
\texttt{sll} & \texttt{ShiftLeftImm} (45), shift & $a = b \cdot 2^{k} \bmod 2^{32}$, $k = \mathit{sa}$ & 4 range \\
\texttt{sllv} & \texttt{ShiftLeft} (62), R & $a = b \cdot 2^{k} \bmod 2^{32}$, $k = c \bmod 32$ & 4 range \\
\texttt{srl}, \texttt{sra}, \texttt{rotr} & \texttt{ShiftRightImm} (83), shift & low word of $(\hat b \gg k)$, $k = \mathit{sa}$ & 1 \texttt{MSB}, 8 \texttt{ShrCarry}, 16 range \\
\texttt{srlv}, \texttt{srav}, \texttt{rotrv} & \texttt{ShiftRight} (89), R & low word of $(\hat b \gg k)$, $k = c \bmod 32$ & 1 \texttt{MSB}, 8 \texttt{ShrCarry}, 16 range \\
\bottomrule
\end{tabular}
\end{center}

A shift by $k = 8q + r$ is a shift by $r$ bits followed by a relabelling of whole bytes by $q$. The shift amount is decomposed into bits, of which the low five are used (the manual's reduction modulo 32), and one-hot selectors for $q \in [0,4)$ and $r \in [0,8)$ are tied to them. A left shift multiplies by $2^{r}$ with witnessed per-byte carries, $\rho_i = 2^{r} b_i + \theta_{i-1} - 256\theta_i$, then moves byte $i$ to byte $i+q$; the shift-amount form computes $2^{r}$ as $(1+s_0)(1+3s_1)(1+15s_2)$ with no one-hot. A right shift first extends $b$ to eight bytes $\hat b$, whose upper four are zero for \texttt{srl}, $\mathtt{0xFF}$ times the sign bit for \texttt{sra}, and a copy of $b$ for \texttt{rotr}, so one right shift leaves the rotation in the low half. The bit part uses eight \texttt{ShrCarry} lookups, each mapping a byte and a 3-bit shift to the shifted byte and the bits shifted out: Jolt's shift decomposition~\citep{jolt}, a subtable indexed by a chunk and the shift amount, with the byte part done by relabelling.

\subsubsection{Branches and jumps}
\label{sec:base-control}

\begin{center}
\small
\begin{tabular}{@{}llll@{}}
\toprule
Instructions & Chip (width), frame & $\mathsf{next\_next\_pc}$ & Added lookups \\
\midrule
\texttt{beq}, \texttt{bne} & \texttt{Branch} (56), I & $\mathsf{next\_pc} + \mathit{off}$ if taken, else $\mathsf{next\_pc} + 4$ & $\le 9$ of 13 \\
\texttt{bgez}, \texttt{bgtz}, \texttt{blez}, \texttt{bltz} & \texttt{Branch} (56), I & as above, comparing with register 0 & $\le 9$ of 13 \\
\texttt{j}, \texttt{jal} & \texttt{Jump} (57), universal & $\mathit{target} \ll 2$ & 3 \texttt{LTU} \\
\texttt{jr}, \texttt{jalr} & \texttt{Jump} (57), universal & $b$ & 3 \texttt{LTU} \\
\texttt{bal} & \texttt{Jump} (57), universal & $\mathsf{next\_pc} + \mathit{off}$ & 3 \texttt{LTU}, 6 range \\
\bottomrule
\end{tabular}
\end{center}

A transfer sets only $\mathsf{next\_next\_pc}$, the second component of the pair it sends; its delay-slot instruction is the next row of the chain (Example~\ref{ex:delay}). \texttt{Branch} compares $a$ (register $\mathsf{rs}$) with $b$ (register $\mathsf{rt}$, or register 0 for the compare-with-zero forms). Equality is two is-zero tests, by witnessed inverses, on the differences of the 16-bit halves; the sign of $a$ is one \texttt{MSB} lookup, gated to the compare-with-zero forms; and $a > 0$ is $(1 - a_s)(1 - [a = 0])$. The offset is shifted and sign-extended at decode, so the taken target is one in-row addition with a byte-checked result, and one \texttt{LTU} lookup per counter word proves it canonical. \texttt{Jump} writes the link $\mathsf{next\_pc} + 4$, which is $\mathsf{pc} + 8$ through the chain, and takes its target from the immediate, from register $b$, or from an addition for \texttt{bal}. As in the executor, \texttt{j} and \texttt{jal} take the region bits as zero (\S\ref{sec:isa-emulator}); the guest toolchain links text below $2^{28}$.

\subsubsection{Loads and stores}
\label{sec:base-memory}

\begin{center}
\small
\begin{tabular}{@{}llll@{}}
\toprule
Instructions & Chip (width), frame & Result & Added interactions \\
\midrule
\texttt{lw}, \texttt{ll} & \texttt{LoadWord} (48), I & $a = M[\mathit{addr}]$ & 9 \\
\texttt{sw}, \texttt{sc} & \texttt{StoreWord} (52), I & $M[\mathit{addr}] \leftarrow a$; \texttt{sc} also sets $a = 1$ & 9 \\
\texttt{lb}, \texttt{lbu}, \texttt{lh}, \texttt{lhu} & \texttt{LoadNarrow} (59), I & selected byte or half, sign- or zero-extended & 10 \\
\texttt{sb}, \texttt{sh} & \texttt{StoreNarrow} (55), I & byte or half of $a$ merged into $M[\mathit{addr}]$ & 9 \\
\texttt{lwl}, \texttt{lwr}, \texttt{swl}, \texttt{swr} & \texttt{MemoryUnaligned} (59), I & byte-wise merge of $M[\mathit{addr}]$ and $a$ & 9 \\
\bottomrule
\end{tabular}
\end{center}

The five memory chips share one block of columns and nine interactions. The effective address $b + \mathrm{sext}(\mathit{off})$ is one in-row addition with range-checked result bytes; one \texttt{LTU} proves it canonical and a second, with an is-zero test on its upper bytes, that it lies above the 36 register addresses; and one \texttt{AND} lookup extracts its two low bits, from which the aligned address is linear. The access is one word-granular access at sub-cycle position 0 (two memory-bus interactions and two timestamp lookups), and the word is not byte-checked again (\S\ref{sec:air-memory}). The word forms assert that the low bits are zero; the narrow and unaligned forms turn them into one-hot offset flags and select or merge bytes polynomially, and a narrow signed load adds one \texttt{MSB} lookup and fills with $\mathtt{0xFF}$ times the bit. \texttt{ll} and \texttt{sc} follow the choices of \S\ref{sec:isa-emulator}.

OpenVM likewise accesses aligned four-cell blocks and selects bytes within the row~\citep[\S3.2.4]{openvm}, where the original Jolt design spends one memory-checking operation per byte~\citep{jolt}; the price is the offset flags and merge columns that make the narrow and unaligned chips the widest of the family.
